\documentclass[11pt]{article}
\usepackage[margin=2cm,headheight=14pt]{geometry}
\usepackage{amsmath,amssymb}
\usepackage{multicol}
\usepackage{fancyhdr}

\pagestyle{fancy}
\fancyhf{}
\fancyhead[L]{Astroestadística}
\fancyhead[R]{Semestre 2026\text{-}2}
\renewcommand{\headrulewidth}{0.4pt}

\begin{document}
\section*{Quiz 1: transformación inversa}

\textbf{Nombre:} \rule{7cm}{0.4pt}\hfill
\textbf{Puntaje:} \rule{1.5cm}{0.4pt} / 10

Cada pregunta vale 1 punto. Marque una sola opción por pregunta.

\begingroup
\setlength{\parskip}{0pt}
\begin{multicols}{2}
\begin{enumerate}
    \item Sea $\widetilde f(x)=x$ para $0\leq x\leq2$ y cero fuera de ese
    intervalo. ¿Cuál es la densidad normalizada?
    \begin{enumerate}
        \item $f(x)=x$ en $[0,2]$.
        \item $f(x)=x/2$ en $[0,2]$.
        \item $f(x)=2x$ en $[0,2]$.
    \end{enumerate}

    \item Una variable tiene densidad $f(x)=3x^2$ para $0\leq x\leq1$.
    ¿Cuál es su función de distribución acumulada en ese intervalo?
    \begin{enumerate}
        \item $F(x)=x^3$.
        \item $F(x)=x^2$.
        \item $F(x)=3x$.
    \end{enumerate}

    \item Para la densidad de la pregunta anterior, se genera $U=1/8$.
    ¿Qué valor produce el método de transformación inversa?
    \begin{enumerate}
        \item $X=1/4$.
        \item $X=1/2$.
        \item $X=3/4$.
    \end{enumerate}

    \item La densidad de $X$ es $f(x)=1/4$ para $0\leq x\leq4$ y cero fuera
    del intervalo. ¿Cuánto vale $F(3)$?
    \begin{enumerate}
        \item $1/4$.
        \item $3/4$.
        \item $1$.
    \end{enumerate}

    \item Una FDA es $F(x)=x/5$ para $0\leq x\leq5$. Si el generador uniforme
    devuelve $U=0.4$, ¿cuál es la observación transformada?
    \begin{enumerate}
        \item $X=0.4$.
        \item $X=2$.
        \item $X=4.6$.
    \end{enumerate}

    \item Una variable discreta toma los valores $-2$, $1$ y $5$ con
    probabilidades $0.25$, $0.50$ y $0.25$, respectivamente. ¿Qué valor se
    obtiene si $U=0.6$?
    \begin{enumerate}
        \item $X=-2$.
        \item $X=1$.
        \item $X=5$.
    \end{enumerate}

    \item Si la FDA tiene tramos constantes, ¿qué definición permite obtener
    correctamente sus cuantiles?
    \begin{enumerate}
        \item $F^{-1}(u)=\inf\{x:F(x)\geq u\}$.
        \item $F^{-1}(u)=1/F(u)$.
        \item $F^{-1}(u)$ es siempre la derivada de $F(u)$.
    \end{enumerate}

    \item Se transforman dos uniformes independientes, $U_1$ y $U_2$, usando
    la misma función inversa. ¿Qué se puede afirmar de $X_1$ y $X_2$?
    \begin{enumerate}
        \item Son independientes.
        \item Siempre toman el mismo valor.
        \item Una observación determina necesariamente la otra.
    \end{enumerate}

    \item La FDA objetivo se puede evaluar, pero no se puede despejar su inversa
    de forma explícita. ¿Qué procedimiento es apropiado para obtener $X$ a
    partir de $U$?
    \begin{enumerate}
        \item Resolver numéricamente $F(x)=U$, por ejemplo, con bisección.
        \item Reemplazar $F(x)$ por la densidad $f(x)$.
        \item Usar $X=F(U)$ sin comprobar la transformación.
    \end{enumerate}

    \item Para una variable generada correctamente por transformación inversa,
    se sabe que $F(4)=0.92$. ¿Qué representa este valor?
    \begin{enumerate}
        \item $\Pr(X=4)=0.92$ necesariamente.
        \item $\Pr(X\leq4)=0.92$.
        \item $\Pr(X>4)=0.92$.
    \end{enumerate}
\end{enumerate}
\end{multicols}
\endgroup
\end{document}